\documentclass[12pt]{article}

\input{../../_assets/preambulo.tex}


\begin{document}

    % 1. Foto de fondo
    % 2. Título
    % 3. Encabezado Izquierdo
    % 4. Color de fondo
    % 5. Coord x del titulo
    % 6. Coord y del titulo
    % 7. Fecha

    
    \input{../../_assets/portada}
    \portadaExamen{ffccA4.jpg}{Ecuaciones\\Diferenciales II\\Examen IV}{Ecuaciones Diferenciales II. Examen IV}{MidnightBlue}{-8}{28}{2026}{David Muñoz Gómez}

    \begin{description}
        \item[Asignatura] Ecuaciones Diferenciales II.
        \item[Curso Académico] 2025/2026.
        \item[Grado] Doble Grado en Informática y Matemáticas.
        \item[Grupo] Único.
        \item[Profesor] Rafael Ortega Ríos.
        \item[Descripción] Convocatoria Extraordinaria.
        \item[Fecha] 9 de Julio de 2026.
        % \item[Duración] ---.
    
    \end{description}
    \newpage


    % ------------------------------------
    
    \begin{ejercicio}
        Construye una retracción del espacio $\bb{R}^3$ en el conjunto $$[0, 1] \times [0, 1] \times [0, 1]$$
    \end{ejercicio}

    \begin{ejercicio}
        Demuestra que la solución maximal del problema
        \[
        x'' + 2x' + x + x^3 = 0, \quad x(0) = 1, \quad x'(0) = 0
        \]
        está definida en un intervalo que contiene a $[0, \infty[$.
    \end{ejercicio}

    \begin{ejercicio}
        Consideramos la sucesión de funciones $\{f_n\}_{n \geq 0}$ definida por
        \[
        f_n : [0, 1] \to \bb{R}, \quad f_n(t) = 2^n \operatorname{sen}(2^{-n} e^t)
        \]
        ¿es uniformemente acotada? ¿es equicontinua?
    \end{ejercicio}

    \begin{ejercicio}
        Dado $\lambda \in \bb{R}$ se considera la ecuación
        \[
        x(t) = \frac{1}{2} x(1) + \lambda \int_0^t x(s) \cos s \, ds, \quad t \in [0, 1],
        \]
        donde $x : [0, 1] \to \bb{R}$ es una función continua.
        \begin{enumerate}[label=\alph*)]
            \item Determina $\lambda_* > 0$ de manera que esta ecuación tenga una única solución si $|\lambda| < \lambda_*$.
            \item Encuentra un número $\lambda^* \in \bb{R}$ para el que la ecuación tiene más de una solución si $\lambda = \lambda^*$.
        \end{enumerate}
    \end{ejercicio}

    \begin{ejercicio}
        Dada la función
        \Func{\varphi}{]0, \infty[}{\bb{R}}{t}{\cos\left(\frac{1}{t}\right)}
        ¿es posible encontrar una ecuación diferencial en el plano de la que $\varphi$ sea solución?
        \begin{observacion}
        Se consideran ecuaciones del tipo $\dot{x} = X(t, x)$ con $X : \bb{R}^2 \to \bb{R}$ continua.
        \end{observacion}
    \end{ejercicio}

    \begin{ejercicio}
        Para cada $\veps \in \bb{R}$ consideramos la ecuación
        \[
        \ddot{x} + 3 \operatorname{sen} x = \veps \operatorname{sen} t.
        \]
        Prueba:
        \begin{enumerate}[label=\alph*)]
            \item Dado $v \in \bb{R}$ existe una única solución $x(t; v, \veps)$ que cumple $x(0) = 0$, $\dot{x}(0) = v$. Además $x(t; v, \veps)$ está definida en todo $\bb{R}$,
            \item la función $F : \bb{R}^2 \to \bb{R}$, $F(v, \veps) = x(\pi; v, \veps)$ es de clase $C^1$,
            \item existe $\veps_* > 0$ de manera que si $|\veps| < \veps_*$ entonces la ecuación tiene una solución que cumple $x(0) = x(\pi) = 0$.
        \end{enumerate}
    \end{ejercicio}

    \newpage
    \setcounter{ejercicio}{0} % Reiniciar contador de ejercicios
    % \noindent
    % \textbf{Solución.}

\end{document}